\documentclass[10pt,a4paper]{amsart}
\usepackage[margin=19mm]{geometry}
\usepackage{amsmath,amssymb,mathtools}
\usepackage[scr=rsfs,cal=euler]{mathalfa}
\usepackage{xfrac}
\usepackage[unicode,hidelinks,bookmarks=false]{hyperref}
\newcommand{\ie}{i.e.\ }
\newcommand{\cf}{cf.\ }
\newcommand{\resp}{respectively\ }
\newcommand{\Subsection}{\subsection}
\providecommand{\nobreakdash}{\nobreak\hbox{-}}
\setlength{\emergencystretch}{2em}
\multlinegap0pt
\usepackage{xcolor}
\usepackage{amssymb}
\usepackage{algorithm}\usepackage{algpseudocode}
\usepackage{mathtools}
\newtheorem{proposition}{Proposition}
\title{Analytic Proof of a Quartic Continued Fraction Identity for $8/\pi^2$ via Operator Decoupling}
\author{Chao Wang}
\date{Source: arXiv:2602.03027v3 (CC BY 4.0)}
\begin{document}
\maketitle

\noindent\emph{Source and licence.} Analytic Proof of a Quartic Continued Fraction Identity for $8/\pi^2$ via Operator Decoupling. Version arXiv:2602.03027v3 (CC BY 4.0), distributed by arXiv under the Creative Commons Attribution 4.0 International licence, http://creativecommons.org/licenses/by/4.0/, as recorded in the arXiv licence metadata for this exact version and shown on the abstract page https://arxiv.org/abs/2602.03027v3. Full licence notice: \texttt{licenses/arxiv-2602.03027-CC-BY-4.0.txt}.\par\smallskip

\noindent\textit{Editorial scope.} Counted unit: the article's proposition prop:minimal (source0.tex lines 297--303). The article's complete original proof of it is delivered with it (source0.tex lines 305--335, one \texttt{proof} environment of 31 lines). No mathematics is written, repaired or re-derived: the statement and the proof are the article's own lines, in the article's own notation, and no retained line is substituted or re-flowed.  No further source-local context is needed: every cross-reference of every retained line resolves inside the two delivered spans.  Nothing was omitted from any retained span, so the package carries no omission record.  No retained line contains a citation, so the wrapper carries no bibliography and no bibliography entry.  No retained line contains a figure environment, an image-inclusion command or drawing code, so no image is delivered and the package declares no asset dependency.  Editorial formatting only, in addition: an AMS \texttt{amsart} preamble that restates the article's own declarations and macros used by the retained lines, namely the environments \texttt{proposition} on separate counters, together with the standard packages those lines need.  The retained environments print in this document as Proposition 1 in order of appearance, whereas the article prints the same items with the numbering of its own full document.  The complete native source of the article as distributed in the arXiv e-print for this exact version is bundled unchanged as \texttt{sources/193878/source0.tex}.
	\begin{proposition}[Structure of the Solution Basis and the Minimal Solution]\label{prop:minimal}
		The numerator sequence $\{A_n\}$ and the denominator sequence $\{B_n\}$ defined in Section~2 constitute a basis of the solution space for the recurrence $\mathcal{L}\,y = 0$.  Both are dominant solutions satisfying $A_n \sim (8/\pi^2)\,B_n$ as $n\to\infty$.  The true minimal solution is the linear combination
		\begin{equation}\label{eq:minimal_def}
			f_n \;=\; A_n - \frac{8}{\pi^2}\,B_n,
		\end{equation}
		which satisfies $\lim_{n\to\infty} f_n/B_n = 0$, confirming that the continued fraction converges.
	\end{proposition}

	\begin{proof}
		Consider the structure of the solutions derived via operator decoupling. From the boundary analysis in Phase 2, the numerator sequence is characterized by the vanishing of its auxiliary component ($w_n \equiv 0$), which forces $\{A_n\}$ to propagate strictly as the pure product trajectory:
		\begin{equation}
			A_n = \prod_{j=1}^{n+1} d_j.
		\end{equation}
		The denominator sequence, conversely, incorporates the propagation of the non-vanishing auxiliary term $w_0 = 1$, leading to the representation:
		\begin{equation}
			B_n = \left( \prod_{j=1}^{n+1} d_j \right) \left( \sum_{k=0}^{n} t_k \right).
		\end{equation}
		To determine the convergence behaviour, we examine the limit of the ratio between these two linearly independent solutions:
		\begin{equation}
			\textcolor{blue}{L \;:=\;} \lim_{n \to \infty} \frac{A_n}{B_n} = \lim_{n \to \infty} \frac{1}{\sum_{k=0}^{n} t_k}.
		\end{equation}
		As established in Phase 3, the general term $t_k$ satisfies the ratio test $\lim_{k \to \infty} |t_{k+1}/t_k| = 1/2 < 1$, which implies the absolute convergence of the series $S = \sum_{k=0}^{\infty} t_k$.
		
		Since $S = \pi^2/8$ is a finite positive constant, we obtain
		\begin{equation}
			L = \frac{1}{S} = \frac{8}{\pi^2},
		\end{equation}
		so $A_n \sim (8/\pi^2)\,B_n$ as $n\to\infty$.  In particular $\lim_{n\to\infty} A_n/B_n = 8/\pi^2 \neq 0$, which means that neither $\{A_n\}$ nor $\{B_n\}$ is the minimal solution; both are dominant solutions of equal asymptotic order.
		The true minimal solution is the linear combination $f_n = A_n - L\,B_n$ introduced in \eqref{eq:minimal_def}.  By linearity of the recurrence, $\{f_n\}$ satisfies $\mathcal{L}\,y=0$.  Moreover,
		\begin{equation}
			\frac{f_n}{B_n} = \frac{A_n}{B_n} - \frac{8}{\pi^2} \;\longrightarrow\; 0 \quad (n \to \infty),
		\end{equation}
		
		so $\{f_n\}$ is subdominant relative to $\{B_n\}$ and hence is the minimal solution in the sense of Pincherle's Theorem.  By that theorem, the existence of $\{f_n\}$ guarantees that the GCF converges, and its value is $\lim_{n\to\infty} A_n/B_n$, which we have already computed:
		\begin{equation}
			\mathcal{K} = S^{-1} = \left( \sum_{m=1}^{\infty} \frac{2^m}{m^2 \binom{2m}{m}} \right)^{-1} = \frac{8}{\pi^2}.
		\end{equation}
		The proof is complete.
	\end{proof}

\end{document}
